

\documentclass[11pt]{article}

\usepackage[margin=1in]{geometry}
\usepackage{amsmath}
\usepackage{amssymb}
\usepackage{bm}
\usepackage{dsfont}
\usepackage[T1]{fontenc}
\usepackage{textcomp}
\usepackage{hyperref}
\usepackage{parskip}
\setcounter{secnumdepth}{2}   % number sections and subsections only
\usepackage{xcolor}

% Flags a point Claude is not sure about; search for \tocheck to find them all.
\newcommand{\tocheck}[1]{\textcolor{red}{[CHECK: #1]}}

\newcommand{\rr}{\mathbf{r}}
\newcommand{\ee}{\hat{\mathbf{e}}}
\newcommand{\KK}{\mathbf{K}}
\newcommand{\LL}{\Lambda}
\newcommand{\re}{\rr^{e}}
\newcommand{\rem}{\rr^{e\prime}}   % mutant equilibrium conformation
\newcommand{\rrij}{\mathbf{r}_{ij}}
\newcommand{\eeij}{\ee_{ij}}
\newcommand{\eeije}{\ee^{e}_{ij}}
\newcommand{\dije}{d^{e}_{ij}}
\newcommand{\dijem}{d^{e\prime}_{ij}}
\newcommand{\lijm}{l'_{ij}}
\newcommand{\kijm}{k'_{ij}}
\newcommand{\dlij}{\delta l_{ij}}
\newcommand{\dkij}{\delta k_{ij}}
\newcommand{\gij}{g_{ij}}
\newcommand{\lij}{l_{ij}}
\newcommand{\kij}{k_{ij}}
\newcommand{\dij}{d_{ij}}

\title{First-order perturbation theory for the generalized elastic network model}
\author{Julian Echave}
\date{2026-08-26}

\makeatletter
\newcommand{\defref}[2]{\@namedef{r@#1}{{#2}{--}}}
% labels below refer to the main report; numbers as in its current version
\defref{sec:changes}{3.2}
\defref{eq:potential}{6}
\defref{eq:equilibrium-condition}{8}
\defref{eq:potential-gradient}{14}
\defref{eq:covariance}{29}
\defref{eq:stress-energy}{48}
\defref{eq:minimum-energy-change}{50}
\defref{eq:entropy-change}{46}
\defref{sec:expansion}{2.2}
\defref{eq:pseudo-inverse}{32}
\makeatother

\begin{document}
\maketitle


\section*{Note}

This is material removed from the main report on the generalized elastic network
model, kept in case it is needed later. Equation references of the form
\texttt{eq:...} that are not defined here refer to that report.

Two results are worth remembering. The change of the Hessian must be projected
onto the internal subspace before any of the eigenvector or covariance formulas
are applied, because the rigid-body directions of the mutant differ from the
reference's at first order; without the projection those formulas carry an error
that does not vanish as the mutation is made smaller. And the change of the
Hessian has two parts of the same order, one from the change of parameters at
fixed conformation and one from the change of conformation at fixed parameters,
the second involving third derivatives of the potential; keeping only the first
is not a consistent first-order approximation.

\section{First-order perturbation theory}
\label{sec:perturbation}

The results of Section~\ref{sec:changes} are exact within the quadratic approximation, but two of them are expensive: $\rem$ requires a minimisation and $\{\lambda'_n\}$ a diagonalisation of $\KK'$.
This section obtains the mutant's structure, normal modes, fluctuations and energies to first order in the size of the mutation, using only quantities of the reference.

\subsection*{The perturbation, its force and its change of the Hessian}

Write the perturbation as a fixed pattern times a scale,
\begin{equation}
\dlij = \sigma\, m_{ij},
\qquad
m_{ij} = 0 \ \text{ unless } i = m \text{ or } j = m,
\label{eq:mutation-scale}
\end{equation}
with the $m_{ij}$ fixed and $\sigma \to 0$ the limit in which the expansion is made; $k$ is assumed differentiable at the reference lengths, so that $\dkij = O(\sigma)$.
The mutant potential is the reference potential plus a perturbation,
\begin{equation}
V'(\rr) = V(\rr;\LL) + W(\rr),
\qquad
W(\rr) \equiv V'(\rr) - V(\rr;\LL),
\label{eq:perturbation}
\end{equation}
and $W$, being a difference of two potentials of the form \eqref{eq:potential}, is a function of the distances alone and is therefore invariant under rigid-body motions.

Two quantities carry the whole of the perturbation's effect.
The first is the force it exerts at the reference equilibrium.
Since the reference gradient vanishes there, by \eqref{eq:equilibrium-condition} and \eqref{eq:potential-gradient},
\begin{equation}
\mathbf{f}_i \equiv -\frac{\partial W}{\partial \rr_i}\bigg|_{\re}
= \sum_{j \ne i} \kijm \big(\dije - \lijm\big) \eeije
= O(\sigma),
\label{eq:mutational-force}
\end{equation}
localised on the mutated site and its neighbours, and with no component along the six rigid-body directions.
The second is the change of the Hessian, $\delta\KK \equiv \KK' - \KK$, also $O(\sigma)$, with $\KK'$ the Hessian of $V'$ at $\rem$ and $\KK$ that of $V$ at $\re$.

Both Hessians are written in a common frame, with $\rem$ superposed on $\re$, and $\delta\KK$ is understood as projected onto the internal subspace, the complement of the rigid-body directions:
\begin{equation}
\delta\KK \;\to\; \mathbf{P}\,\delta\KK\,\mathbf{P},
\qquad
\mathbf{P} = \KK^{+}\KK = \sum_{\lambda_n \ne 0} \mathbf{q}_n\mathbf{q}_n^{\mathsf T}.
\label{eq:internal-projection}
\end{equation}
This is not a convention but a requirement.
Both networks have exactly six null directions, but they do not span the same subspace: the three translations are the same vectors in both, since a uniform displacement does not depend on the structure, whereas the three infinitesimal rotations are generated by the node positions and so differ at $O(\sigma)$.
An unprojected $\delta\KK$ therefore carries blocks coupling the internal and rigid-body subspaces, and these spoil the formulas below at the order at which they are meant to hold.

\subsection*{The change in the equilibrium conformation}

The mutant equilibrium satisfies $\nabla V'(\rem) = \mathbf{0}$ exactly.
Writing $\rem = \re + \mathbf{w}$ and expanding about $\re$,
\begin{equation}
\mathbf{0} = -\mathbf{f} + \big(\KK + \delta\KK\big)\mathbf{w} + O(\mathbf{w}^{2}),
\end{equation}
in which $\delta\KK\,\mathbf{w}$ is $O(\sigma^{2})$ because $\mathbf{f}$ and $\delta\KK$ are both $O(\sigma)$, so $\KK$ may replace $\KK'$ at this order:
\begin{equation}
\boxed{
\mathbf{w} \equiv \rem - \re = \KK^{+}\mathbf{f} + O(\sigma^{2}) .
}
\label{eq:first-order-displacement}
\end{equation}
This is the linear response approximation: the mutant structure is the reference deformed by the force \eqref{eq:mutational-force}, with $\KK^{+}$ as its compliance.
In the normal-mode representation, writing $w_n = \mathbf{q}_n^{\mathsf T}\mathbf{w}$ for the components of this particular displacement,
\begin{equation}
\boxed{
w_n = \frac{\mathbf{q}_n^{\mathsf T}\mathbf{f}}{\lambda_n} + O(\sigma^{2}),
\qquad \lambda_n \ne 0,
}
\label{eq:first-order-displacement-modal}
\end{equation}
so each mode is excited in proportion to its overlap with the force and in inverse proportion to its force constant.
A localised force gives a delocalised displacement, carried mostly by the modes that are both soft and aligned with it.

\subsection*{The change in the normal modes}

The eigenvalues shift by the diagonal elements of $\delta\KK$ in the reference basis,
\begin{equation}
\boxed{
\lambda'_n = \lambda_n + \mathbf{q}_n^{\mathsf T}\, \delta\KK\, \mathbf{q}_n + O(\sigma^{2}),
}
\label{eq:eigenvalue-shift}
\end{equation}
and the eigenvectors rotate by the off-diagonal ones,
\begin{equation}
\boxed{
\mathbf{q}'_n = \mathbf{q}_n + \sum_{\substack{m \ne n \\ \lambda_m \ne 0}}
\frac{\mathbf{q}_m^{\mathsf T}\, \delta\KK\, \mathbf{q}_n}{\lambda_n - \lambda_m}\, \mathbf{q}_m
+ O(\sigma^{2}) .
}
\label{eq:eigenvector-shift}
\end{equation}
The eigenvalue shifts need no non-degeneracy: when several eigenvalues coincide, their individual shifts are ambiguous but their sum is the trace of $\delta\KK$ on the degenerate subspace, which is not.
The eigenvector shifts do: \eqref{eq:eigenvector-shift} requires the gaps $\lambda_n - \lambda_m$ to be large compared with the corresponding elements of $\delta\KK$, and two modes of nearly equal force constant mix strongly under any perturbation that couples them.
Since $V'$ is invariant under rigid-body motions like $V$, the mutant also has exactly six null eigenvalues, and \eqref{eq:eigenvalue-shift} returns zero for them; the directions they span, however, are the mutant's own, as noted above.

\subsection*{The change in the fluctuations}

Differentiating $\bm{\Sigma} = \beta^{-1}\KK^{+}$ of \eqref{eq:covariance}, and using \eqref{eq:internal-projection},
\begin{equation}
\boxed{
\delta\bm{\Sigma} = -\beta^{-1}\, \KK^{+}\, \delta\KK\, \KK^{+} + O(\sigma^{2}) ,
}
\label{eq:covariance-change}
\end{equation}
which requires no diagonalisation of either network.
In the reference normal-mode basis its elements are
\begin{equation}
\boxed{
\mathbf{q}_m^{\mathsf T}\, \delta\bm{\Sigma}\, \mathbf{q}_n
= -\frac{1}{\beta}\, \frac{\mathbf{q}_m^{\mathsf T}\, \delta\KK\, \mathbf{q}_n}{\lambda_m \lambda_n}
+ O(\sigma^{2}) ,
}
\label{eq:covariance-change-modal}
\end{equation}
whose diagonal, $\delta\langle u_n^{2}\rangle = -\delta\lambda_n/\beta\lambda_n^{2}$, is the change in each mode's amplitude: stiffening a mode narrows its distribution, and the weight $\lambda_n^{-2}$ makes the change largest for the softest modes.
The off-diagonal elements are the substance of the change of basis: in the mutant the reference modes are no longer uncorrelated, and the correlation between two of them is set by the element of $\delta\KK$ that couples them, divided by both force constants.

\subsection*{The changes in the energies}

The stress energy \eqref{eq:stress-energy} needs no expansion: it is already exact and evaluated at the reference.
The relaxation energy follows from \eqref{eq:first-order-displacement},
\begin{equation}
\boxed{
\Delta V_{\text{relax}} = -\tfrac12\, \mathbf{f}^{\mathsf T}\KK^{+}\mathbf{f} + O(\sigma^{3})
= -\frac{1}{2}\sum_{\lambda_n \ne 0} \frac{\big(\mathbf{q}_n^{\mathsf T}\mathbf{f}\big)^{2}}{\lambda_n}
+ O(\sigma^{3}) ,
}
\label{eq:first-order-relaxation}
\end{equation}
non-positive, as the exact result requires, and $O(\sigma^{2})$: the structural response reaches the energy one order later than the stress.
With \eqref{eq:minimum-energy-change} this gives $\Delta U = \Delta V_{\min}$ without the mutant structure.

For the entropy, substituting \eqref{eq:eigenvalue-shift} into \eqref{eq:entropy-change} and expanding the logarithm,
\begin{equation}
\boxed{
T\Delta S = -\frac{1}{2\beta} \sum_{\lambda_n \ne 0} \frac{\delta\lambda_n}{\lambda_n}
= -\frac{1}{2\beta} \operatorname{tr}\big(\KK^{+} \delta\KK\big) + O(\sigma^{2}) ,
}
\label{eq:first-order-entropy-change}
\end{equation}
which is $O(\sigma)$ and needs neither spectrum: only $\KK^{+}$ and $\delta\KK$ enter, and the trace projects onto the internal subspace of its own accord.
A mutation that stiffens the network lowers its entropy.
The free energy change is then
\begin{equation}
\boxed{
\Delta A = \Delta V_{\text{stress}} - \tfrac12 \mathbf{f}^{\mathsf T}\KK^{+}\mathbf{f}
+ \frac{1}{2\beta} \operatorname{tr}\big(\KK^{+} \delta\KK\big) + O(\sigma^{3}) .
}
\label{eq:first-order-free-energy-change}
\end{equation}

\subsection*{What is assumed}

Four things, none of them required by Section~\ref{sec:changes}.

$k$ is differentiable at the reference lengths, so that $\dkij = O(\sigma)$.
The exact results hold for any non-negative $k$.

The displacement \eqref{eq:first-order-displacement} treats $\KK$ as constant between $\re$ and $\rem$, which is consistent at this order but fails once the mutation moves the network appreciably along its soft directions.

The eigenvector shifts \eqref{eq:eigenvector-shift} and the expansions of the logarithm require the perturbation to be small compared with the spectral gaps and with the $\lambda_n$ themselves, which the softest modes violate first.

The change of the Hessian has two parts of the same order,
\begin{equation}
\delta\KK =
\underbrace{\frac{\partial^2 W}{\partial \rr^2}\bigg|_{\re}}_{\text{parameters}}
\;+\;
\underbrace{\frac{\partial^3 V}{\partial \rr^3}\bigg|_{\re} \cdot \big(\KK^{+}\mathbf{f}\big)}_{\text{structure}}
\;+\; O(\sigma^{2}),
\label{eq:hessian-change}
\end{equation}
the second involving third derivatives of $V$, outside the quadratic approximation of Section~\ref{sec:expansion}.
Dropping it and keeping only the change of parameters at fixed conformation is not a consistent approximation: the two are of the same order in $\sigma$, so the error does not vanish as the mutation is made smaller.


\end{document}
